\section{The Finite Certificate}\label{sec:certificate}

In this section we fix the finite data underlying Theorem \ref{thm:main}, namely the set $T$, the selected set $S$, and the weights $k(p)$, and we carry out the verification of the criterion (\ref{eq:c3}) and the evaluation of the master inequality.  Throughout, $T$ denotes the set of the first $38$ odd primes and $D=\prod_{q\in T}q$.  We found the data by a computer search over candidate sets $T$ of between $30$ and $50$ small odd primes, combined, for each candidate, with an exact optimization of $S$, the weights, and the radius $R$.  Our first goal will be to certify the analytic input $\log \mathrm{B}_{E_{39}}(\lambda)<0.09$ used in the master inequality.

\subsection{The $E_{39}$ Louboutin Certificate}\label{sec:e39}

Put
\[
   E_{39}=\mathbb{Q}(\sqrt q:q\in T,\sqrt2),
\]
a multiquadratic field of degree $2^{39}$.  Proposition \ref{thm:fields} guarantees that every tower field $F$ contains $E_{39}$, and so we may apply Theorem \ref{prop:master} with $E=E_{39}$.  The point of using the full multiquadratic field $E_{39}$, rather than a smaller subfield, is that each ramified prime $q\in T$ and the prime $2$ then contribute negative correction terms in the expression (\ref{eq:euler39}) for the degree normalized $\log\zeta_{E_{39}}$.

We begin with the local data of $E_{39}$:
\begin{lemma}\label{lem:E39local}
The field $E_{39}$ has degree $2^{39}$, and its local factors have the following form.  At $p=2$ the completion is $\mathbb{Q}_2(\sqrt3,\sqrt5,\sqrt2)$, with $e=4$ and $f=2$.  At each $q\in T$ we have $e=2$ and $f=2$, where the residual quadratic direction is witnessed by some $q'\in T$ with $\left(\frac{q'}{q}\right)=-1$, listed for every $q$ in the second table of the Appendix.  At every other prime $p$ the Frobenius has order at most $2$, and $p$ splits completely in $E_{39}$ if and only if $\left(\frac{2}{p}\right)=1$ and $\left(\frac{q}{p}\right)=1$ for all $q\in T$.  Consequently, for $s>1$,
\begin{equation}\label{eq:euler39}
\begin{aligned}
   \log\zeta_{E_{39}}(s)^{1/2^{39}}
   &=\frac12\log\zeta(2s)
   -\frac14\sum_{q\in T}\left(-\log(1-q^{-2s})\right)  \\
   &\quad -\frac38\left(-\log(1-2^{-2s})\right)
   +R_{39}(s),
\end{aligned}
\end{equation}
where
\[
   R_{39}(s)=\sum_{p\text{ split completely in }E_{39}}\artanh(p^{-s}).
\]
\end{lemma}

\begin{proof}
We first treat the dyadic place.  The square classes of odd rational primes in $\mathbb{Q}_2^\times/\mathbb{Q}_2^{\times2}$ are determined by their residues modulo $8$ and are generated by the classes of $3$ and $5$.  Since $3,5\in T$, the dyadic completion is $\mathbb{Q}_2(\sqrt3,\sqrt5,\sqrt2)$.  The class $5$ gives the unramified quadratic direction and the classes $3$ and $2$ give total ramification index $4$, so that $e=4$ and $f=2$.

For an odd ramified prime $q\in T$, ramification comes only from $\sqrt q$, and any witness $q'\in T$ with $\left(\frac{q'}{q}\right)=-1$ forces a residual unramified quadratic direction, so $e=f=2$.  At primes outside $2T$ the multiquadratic extension is unramified with elementary abelian Galois group, so Frobenius has order $1$ or $2$, and complete splitting is exactly the simultaneous vanishing of all the quadratic characters, namely $\left(\frac{2}{p}\right)=\left(\frac{q}{p}\right)=1$ for all $q\in T$.

We now compute the degree normalized Euler factors.  A generic unramified prime with Frobenius of order $2$ contributes local exponent $\frac1{ef}=\frac12$, giving the factor $(1-p^{-2s})^{-1/2}$ and in total the term $\frac12\log\zeta(2s)$ up to the corrections below.  A completely split prime contributes $(1-p^{-s})^{-1}$ instead, and
\[
   \frac{(1-p^{-s})^{-1}}{(1-p^{-2s})^{-1/2}}
   =\exp(\artanh(p^{-s})),
\]
which accounts for $R_{39}(s)$.  A ramified odd $q\in T$ has normalized local exponent $\frac1{ef}=\frac14$, a deficit of $\frac14$ against the generic $\frac12$, and the dyadic exponent is $\frac1{ef}=\frac18$, a deficit of $\frac38$.  This yields (\ref{eq:euler39}).
\end{proof}

Armed with Lemma \ref{lem:E39local}, we are now ready to prove the required bound:

\begin{proposition}\label{prop:B39}
With $x_0=3\cdot10^{-11}$, $s_0=1+x_0$, and $\lambda=\sqrt{16D}$, we have that
\[
   \log \mathrm{B}_{E_{39}}(\lambda)<0.09.
\]
\end{proposition}

\begin{proof}
We evaluate the infimum (\ref{eq:B-def}) at the single point $x=x_0$ and bound the five resulting pieces.  Each numerical bound below is a strict upper bound obtained from evaluation with interval arithmetic.

First, a high precision evaluation of the explicit part of (\ref{eq:euler39}) gives
\[
   \frac12\log\zeta(2s_0)
   -\frac14\sum_{q\in T}\left(-\log(1-q^{-2s_0})\right)
   -\frac38\left(-\log(1-2^{-2s_0})\right)
   <0.0887019.
\]

Second, we bound the initial segment of $R_{39}(s_0)$.  Complete splitting in $E_{39}$ implies complete splitting in the subfield $E_{30}$ generated by the square roots of the first thirty odd primes, so a deterministic enumeration of the primes up to $10^{13}$ that split completely in $E_{30}$ gives
\[
   \sum_{\begin{array}{c}
      \scriptstyle p\leq10^{13}\\
      \scriptstyle p\text{ split completely in }E_{39}
   \end{array}}\artanh(p^{-s_0})
   \leq
   \sum_{\begin{array}{c}
      \scriptstyle p\leq10^{13}\\
      \scriptstyle p\text{ split completely in }E_{30}
   \end{array}}\artanh(p^{-s_0})
   <7.34\cdot10^{-11}.
\]

Third, we cover the range from $10^{13}$ to $8D$ by a Selberg upper bound sieve \cite{HalberstamRichert,IwaniecKowalski} in dyadic blocks, applied within the two admissible residue classes $p\equiv\pm1\pmod{8}$.  If $C(Y)$ denotes the number of completely split primes in $(Y,2Y]$, then for every $Q\geq1$ we have that
\[
   C(Y)\leq \frac{2\,(Y/8+1+Q^2)}{H_{38}^*(Q)},
\]
where
\[
   H_{38}^*(Q)=\sum_{\begin{array}{c}
      \scriptstyle d\mid\prod_{q\in T}q\\
      \scriptstyle d\leq Q
   \end{array}}\ \prod_{q\mid d}\frac{q+1}{q-1}.
\]
With $Y_j=10^{13}\cdot2^j$, $Q_j=\lfloor0.315\lfloor\sqrt{Y_j}\rfloor\rfloor$, and $0\leq j\leq178$, an exact meet-in-the-middle evaluation of the sums $H_{38}^*(Q_j)$ gives
\[
   \sum_{j=0}^{178}
   \frac{2\,(Y_j/8+1+Q_j^2)}{H_{38}^*(Q_j)}\,
   \artanh(Y_j^{-s_0})<0.00006066.
\]

Fourth, for the tail $Y_j>8D$ the Brun-Titchmarsh theorem of Montgomery and Vaughan \cite{MontgomeryVaughan} gives
\[
   \sum_{j\geq179} C(Y_j)\artanh(Y_j^{-s_0})<2\cdot10^{-10}.
\]

Fifth, the archimedean factor satisfies
\[
   \frac{x_0}{2}(\log\lambda-\log\pi)+\log\Gamma(1+x_0/2)<1.2\cdot10^{-9}.
\]

Adding the five bounds, we obtain
\[
   \log \mathrm{B}_{E_{39}}(\lambda)
   <0.0887019+7.34\cdot10^{-11}+0.00006066+2\cdot10^{-10}+1.2\cdot10^{-9}
   <0.09.
   \qedhere
\]
\end{proof}

\subsection{The Certificate Data}

We now turn to the discrete data of the certificate.  Every discrete assertion below is verified by a finite computation in exact integer or rational arithmetic, and a verification script accompanies the source of this paper.  Explicitly,
\[
\begin{aligned}
T=\{&3,5,7,11,13,17,19,23,29,31,37,41,43,47,53,59,61,67,71,73,\\
&79,83,89,97,101,103,107,109,113,127,131,137,139,149,151,157,163,167\},
\end{aligned}
\]
the set of the first $38$ odd primes.  It satisfies
\[
   |T|=38,\ \ \ D=\prod_{q\in T}q\equiv3\pmod{8},\ \ \ \text{and}\ \ \ \#\left\{q\in T:\ q\equiv3\pmod{4}\right\}=21,
\]
so the hypotheses of section \ref{sec:dyadic} hold (including $7\in T$), and $g=|T|=38$.

For an odd prime $p\notin T$, we call $p$ \emph{split} or \emph{inert} according as $\left(\frac{D}{p}\right)=1$ or $-1$, and we call an odd prime \emph{admissible} if it satisfies (\ref{eq:admissible}).  We define $S$ to consist of the prime $2$, every odd prime $p\leq283$, every admissible split prime $p\in(283,1973]$, and every admissible inert prime $p\in(283,2267]$.  For this $T$, every odd prime in these two intervals is admissible.  The largest selected prime is $2267$.

The weights are
\[
\begin{array}{c|ccccccccc}
 p&2&3&5&7&11&13,17,19&23\leq p\leq53&59\leq p\leq283&\text{otherwise}\\
\hline
 k(p)&19&12&8&6&5&4&3&2&1,
\end{array}
\]
where the ranges refer to primes of $S$ in those ranges.  The complete table, including splitting status and admissibility witnesses, is given in the Appendix.

We summarize the discrete half of the certificate:
\begin{proposition}\label{prop:t38}
Let $T$ be the set of the first $38$ odd primes, and let $S$ and the weights $k(p)$ be as defined above.  Then
\[
   |S|=314\ \ \ \text{and}\ \ \ U_S=114,
\]
and the weighted Golod-Shafarevich polynomial of Theorem \ref{thm:dyadic-gs},
\[
   P(\tau)=1-38\tau+(38+314+3)\tau^2+114\tau^3
   =1-38\tau+355\tau^2+114\tau^3,
\]
satisfies
\[
   P\!\left(\frac6{115}\right)=-\frac{101}{1520875}<0.
\]
In particular the criterion (\ref{eq:c3}) holds for $(T,S,k)$.
\end{proposition}

\begin{proof}
Each assertion is a finite computation in exact integer or rational arithmetic from the definitions above.  The type counts are
\[
\begin{array}{c|cccc}
\text{type} & \text{dyadic} & \text{ramified} & \text{inert} & \text{split}\\
\hline
\text{count} & 1&38&161&114
\end{array}
\]
and the weight histogram is
\[
\begin{array}{c|ccccccccc}
 k&1&2&3&4&5&6&8&12&19\\
\hline
 \#\{p\in S:k(p)=k\}&253&45&8&3&1&1&1&1&1.
\end{array}
\]
Every odd selected prime is either $1\pmod{4}$ or carries the inertness witness $q\in T$ displayed in the Appendix.  The minimum of $P$ on $(0,1)$ occurs at
\[
   \tau=0.05220818647952869\ldots\ \ \ \text{with}\ \ \ P(\tau)=-6.68471161139\cdot10^{-5}\ldots,
\]
and $\frac6{115}$ is a nearby exact rational witness, which gives the displayed value.
\end{proof}

\noindent Theorem \ref{thm:dyadic-gs} and Proposition \ref{thm:fields} thus supply the required family of fields for Theorem \ref{prop:master}.

\subsection{The Master Inequality Evaluation}

To prove Theorem \ref{thm:main}, what remains is to evaluate the master inequality.
\begin{proof}[Proof of Theorem \ref{thm:main}]
The criterion (\ref{eq:c3}) holds by Proposition \ref{prop:t38}, and $7\in T$, and so Proposition \ref{thm:fields} supplies totally real Galois extensions $F/\mathbb{Q}$ of arbitrarily large degree containing $E_{39}$ such that the CM fields $K=F(i)$ are Galois, every prime of $F$ above every $p\in S$ splits in $K/F$, the discriminant parameters are $(\Delta_K/\Delta_F)^{1/d}=\lambda=\sqrt{16D}$ and $\sqrt{|\Delta_K|}=\lambda^d$, and the local parameters $e(p)$ and $f(p)\leq2$ are as displayed there.  That is, the hypotheses of Theorem \ref{prop:master} hold with $E=E_{39}$ and $\lambda=\sqrt{16D}$.

For the selected prime weights, using $f(p)=2$ for all $p\in S$,
\begin{equation}\label{eq:sigma-sum}
   \Sigma:=\sum_{p\in S}\frac{\log(k(p)+1)}{4e(p)}
   =56.2819843775269\ldots,
\end{equation}
and
\begin{equation}\label{eq:logA-sum}
   \log A=\sum_{p\in S}\frac{k(p)\log p}{2e(p)}
   =1093.4272954355824\ldots .
\end{equation}
The dyadic root discriminant parameter is
\[
   \lambda=\sqrt{16D},\ \ \ \text{so that}\ \ \ \log(\lambda^2)=154.01230121116885\ldots.
\]
Any $R>1/2$ gives a valid bound in Theorem \ref{prop:master}, and we choose the numerically optimal
\[
   R=2.688292831485591\ldots.
\]
Inserting the rigorous bound $\log\mathrm{B}_{E_{39}}(\lambda)<0.09$ of Proposition \ref{prop:B39} into (\ref{eq:N-master}) and (\ref{eq:M-master}) yields
\[
   N(R)>37.8867825632423\ \ \ \text{and}\ \ \ M(R)<1057.332065023008.
\]
Here
\[
   \log\left(\frac{\lambda}{2A}\right)=-1017.114292\ldots,
\]
so the thickening term $\lambda/(2A)$ is negligible at the displayed precision, although we evaluate the exact expression in (\ref{eq:M-master}).  Consequently
\[
   \frac{N(R)}{M(R)}>0.0358324.
\]
Taking $\varepsilon=N(R)/M(R)-0.0358324>0$ in Theorem \ref{prop:master} gives arbitrarily large finite sets $U\subset\mathbb{R}^2$ with at least $(\#U)^{1.0358324}$ ordered unit distance pairs, which proves the theorem.
\end{proof}
